% ============================================================================
% Chapter 6 — An Empirical Program for Testing DLVT
%
% Materializes the validation roadmap of §\ref{sec:roadmap} as a concrete
% measurement model, design, estimation strategy, identification plan, and
% pre-registered killer test (draft pre-registration: prereg/PREREGISTRATION.md
% at the repository root; the parameter-recovery prerequisite is executable in
% the companion module dlvt/recovery.py).
%
% All \citep/\citet keys below are verified against ../references.bib,
% including the methodology references added in v2.1 (Podsakoff 2003,
% Vandenberg & Lance 2000, Lindell & Whitney 2001, Reise 2012, Sliwinski
% 2008, Bell 1953, Beal 2015, Voelkle et al. 2012, Driver et al. 2017,
% Carpenter et al. 2017, Lakens 2017, Rizopoulos 2012, Hahn et al. 2011,
% Bracken et al. 2001).
% ============================================================================
\section{An Empirical Program for Testing DLVT}
\label{sec:empirical}

The propositions of \S\ref{sec:propositions} and the roadmap of
\S\ref{sec:roadmap} declare \emph{what} would falsify the theory. This
section specifies \emph{how}: a measurement model mapping the four model
quantities onto validated instruments (\S\ref{sec:emp-measurement}), a
longitudinal design whose cadence is derived from the verified timescales of
the baseline system (\S\ref{sec:emp-design}), an estimation strategy that
targets the identifiable dimensionless groups rather than the raw parameter
vector (\S\ref{sec:emp-estimation}), an identification plan that breaks the
$\mu/\alpha$ and $R/\delta$ ridges with exogenous variation
(\S\ref{sec:emp-identification}), the pre-registered field experiment that
could kill the theory (\S\ref{sec:emp-killer}), and the threats and ethical
constraints that bound the whole program (\S\ref{sec:emp-threats}). A draft
pre-registration of the field experiment accompanies the repository
(\texttt{prereg/PREREGISTRATION.md}).

\subsection{Measurement model: from constructs to instruments}
\label{sec:emp-measurement}

The model quantities $V$, $C$, $O$, and $I$ are theoretical states, not
observables; each requires an explicit auxiliary measurement theory.
Table~\ref{tab:measurement} summarizes the proposed mapping; the paragraphs
below record the substantive commitments and the places where measurement
choices could themselves be wrong.

\begin{table}[t]
\centering
\caption{Measurement model: construct $\to$ instrument $\to$ indicator and
cadence. ESM = experience-sampling (daily burst) measurement; wave =
quarterly panel wave (\S\ref{sec:emp-design}).}
\label{tab:measurement}
\vspace{0.4em}
\small
\begin{tabular}{p{0.16\linewidth}p{0.36\linewidth}p{0.38\linewidth}}
\toprule
Construct & Instrument & Indicators and cadence \\
\midrule
Vitality $V$ &
Subjective Vitality Scale, state version
\citep{ryan1997subjective,nix1999revitalization}; Recovery Experience
Questionnaire \citep{sonnentag2007recovery} for recovery-side mediators;
MBI emotional-exhaustion subscale \citep{maslach1981burnout} as
discriminant marker &
SVS in every daily ESM burst; REQ (detachment, relaxation, control,
mastery) daily in bursts; MBI-exhaustion at each quarterly wave \\
\addlinespace
Career capital $C$ &
Two-factor status/competence composite: \emph{status}---role tenure and
level, span of control, budget authority, advice-network centrality,
external reputation markers \citep{becker1993human,cross2003energy};
\emph{competence}---structured assessment-center or validated
managerial-skill ratings &
Archival HR and network indicators at each quarterly wave; competence
assessment annually; rival bifactor CFA pre-registered (see text) \\
\addlinespace
Complexity $O$ &
Executive job demands \citep{hambrick2005executive,zhu2022executive};
Work Design Questionnaire complexity and information-processing subscales
\citep{morgeson2006wdq,parker2017workdesign}; structural load metrics
\citep{garicano2000hierarchies} &
Perceptual scales at each quarterly wave; structural metrics (direct
interfaces, meeting hours from calendar metadata, matrixed reporting
lines, decision-queue depth) computed continuously, aggregated
quarterly \\
\addlinespace
Impact $I$ &
Multi-source (360-degree) leadership-effectiveness ratings from
supervisor, peer, and subordinate panels \citep{bracken2001multisource};
objective unit performance (unit P\&L or KPI attainment, unit engagement)
\citep{hambrick2007upper,edmans2016executive} &
360 ratings semi-annually; objective unit indicators quarterly from
archival systems, never self-reported \\
\bottomrule
\end{tabular}
\end{table}

\paragraph{Vitality.}
$V$ is operationalized as state subjective vitality
\citep{ryan1997subjective,ryan2000self}, the construct the model was
micro-founded on, with the Recovery Experience Questionnaire
\citep{sonnentag2007recovery} measuring the recovery-side processes that the
parameter $R$ summarizes. Because a low-vitality equilibrium is not the same
claim as clinical burnout, the MBI emotional-exhaustion subscale
\citep{maslach1981burnout,maslach2001job} is administered as a
\emph{discriminant} marker: vitality and exhaustion should be strongly
negatively correlated but separable in CFA
\citep{halbesleben2004burnout,alarcon2011meta}; if they collapse onto a
single factor, the model's $V$ is not empirically distinct from reversed
exhaustion and the measurement theory must be revised before any dynamic
claim is tested.

\paragraph{Career capital and the scalar-$C$ problem.}
The model treats $C$ as a single stock. Empirically this conflates at least
two components that need not move together: \emph{status} capital (tenure,
span, centrality, reputation---the assets that generate complexity load) and
\emph{assessed competence} (the assets that generate impact). The
distinction matters because promotion systems can raise status while
competence fit falls \citep{pluchino2010peter}. We therefore commit, in
advance, to a rival measurement model: a bifactor CFA
\citep{reise2012bifactor} with a general capital factor plus status and
competence group factors, estimated at every wave. If the general factor is weak (e.g., explained common variance below
a pre-declared threshold) or the two group factors show divergent
trajectories, the scalar-$C$ specification is rejected \emph{as
measurement}, and the dynamic model must be re-derived with a two-stock
capital vector before Propositions P1--P6 can be tested. This is a
deliberate exposure: the theory should not be allowed to survive by
absorbing a measurement failure into parameter estimates.

\paragraph{Complexity and impact.}
$O$ combines the perceptual executive-job-demands tradition
\citep{hambrick2005executive,zhu2022executive} and the complexity and
information-processing subscales of the Work Design Questionnaire
\citep{morgeson2006wdq,parker2017workdesign} with structural metrics
(interfaces, coordination hours, reporting lines) in the spirit of
\citet{garicano2000hierarchies}; the challenge--hindrance distinction
\citep{podsakoff2023challenge} is recorded so that the drain kernel can be
probed for demand-type heterogeneity (P4). $I$ is measured from sources
other than the focal leader: multi-source ratings and objective unit
performance \citep{hambrick2007upper,edmans2016executive}.

\paragraph{Longitudinal invariance and common-method defenses.}
Two requirements apply to every construct. First, \emph{longitudinal
measurement invariance} \citep{vandenberg2000invariance}: configural,
metric, and scalar invariance across waves must be established (or partial
invariance explicitly modeled) before latent trajectories are interpreted
as state dynamics; without scalar invariance, an apparent damped spiral can
be manufactured by drifting item intercepts. Second,
\emph{common-method-bias defenses} \citep{podsakoff2003common}: the design
separates sources ($V$ self-reported; $I$ other-reported and archival; $C$
and the structural half of $O$ archival), separates measurement occasions
(ESM bursts versus quarterly waves), and includes a theoretically neutral
marker variable \citep{lindell2001marker} in every wave so that method
variance can be partialled in sensitivity analyses. Because the central predictions are
\emph{within-person dynamic signatures}---decay rates, phase relations,
conservation of $\beta C^{\eta}$---rather than cross-sectional correlations,
stable individual response styles difference out; but occasion-specific
method variance does not, and the defenses above are treated as mandatory,
not optional.

\subsection{Design: measurement bursts inside an accelerated panel}
\label{sec:emp-design}

The verified linearization at the baseline equilibrium has eigenvalues
$-0.205\pm0.331i$ (Table~\ref{tab:parameters}), i.e.\ a damped spiral with
relaxation time $1/0.205\approx4.9$ time units and rotation period
$2\pi/0.331\approx19$ time units, while the fast vitality timescale is
$1/R\approx0.33$ time units against the slow capital timescale
$1/\mu=5$---a separation of roughly $15\times$ in raw rates (the
\emph{eigenmode} separation at the equilibrium is milder, $\approx3\times$,
with mixed complex modes---Appendix~\ref{app:fastslow}---which is why the
design does not rely on a strict fast--slow split). Under the illustrative
weekly calibration (\S\ref{sec:calibration}) these translate to a vitality
boundary layer of about two working days, capital relaxation of about five
weeks, and a spiral period of about $19$ weeks. No single sampling cadence
can resolve both ends of this timescale spread economically. The
design therefore has two interlocking clocks:

\begin{itemize}
\item \emph{Daily ESM bursts for the fast $V$ dynamics}
\citep{beal2015esm,sliwinski2008burst}. Two working weeks
(10 working days) of daily short-form measurement (state SVS, REQ items,
momentary load) once per quarter. Daily sampling oversamples the
$1/R\approx0.33$-unit boundary layer by design and, by Nyquist logic,
resolves any within-burst oscillation with period longer than two days;
its purpose is to estimate the fast recovery--drain balance
(quasi-equilibrium $V_{\mathrm{qe}}(O)$ of Proposition~\ref{prop:band}(iii))
and the diffusion scale of $V$.
\item \emph{Quarterly panel waves for the slow $C$, $O$, $I$ dynamics.}
Twelve to thirteen waves over three years carrying the full
Table~\ref{tab:measurement} battery. Sampling the $19$-unit spiral at
face value requires an interval below the Nyquist bound of $9.5$ time
units ($\approx9.5$ weeks under the weekly calibration); a rigid
$13$-week quarterly grid sits above that bound. Three features close the
gap: (i) wave timing is deliberately \emph{staggered across participants}
(uniformly jittered within each quarter), so the pooled design is
irregularly sampled and does not alias a common frequency; (ii) the
estimator is continuous-time (\S\ref{sec:emp-estimation}) and uses exact
inter-observation intervals rather than assuming a uniform lag; and
(iii) the embedded bursts pin the fast dynamics locally, constraining the
phase of the slow spiral between waves. If pilot bursts indicate real
timescales as fast as the illustrative calibration, brief monthly pulse
surveys (three items) are added; the calibration is illustrative, and the
mapping from model time to calendar time is itself an estimand, so cadence
is re-derived from pilot estimates rather than trusted a priori.
\end{itemize}

\paragraph{Accelerated longitudinal structure.}
Because equilibrium claims concern career horizons far longer than three
years, the panel is \emph{accelerated} \citep{bell1953convergence}:
participants are enrolled in staggered role-tenure cohorts (roughly $0$--$2$, $3$--$5$, $6$--$10$, and
$>\!10$ years in comparable executive roles), so that three calendar years
of observation span two decades of career time under a convergence
assumption. The convergence assumption---that cohorts trace segments of a
common trajectory---is testable via cohort-by-time interactions and is
itself informative for P1: a single-attractor theory predicts that
late-tenure cohorts cluster near $(V^{*},C^{*})$ with small oscillations,
while early-tenure cohorts display the large-amplitude transient
(overshoot--crash--recovery) of the damped spiral. Persistent
between-cohort divergence not explained by parameters would count against
the single-attractor claim (P1's falsifier).

\subsection{Estimation: continuous time, reduced targets}
\label{sec:emp-estimation}

\paragraph{Continuous-time state-space form.}
Discrete-time cross-lagged models estimate lag-specific coefficients that
change with the sampling interval and cannot express the model's actual
claim, which is a drift field. Estimation therefore uses continuous-time
structural equation modeling
\citep{voelkle2012ctsem,driver2017ctsem} or a Bayesian
state-space model in Stan \citep{carpenter2017stan} whose drift is the
\emph{exact} DLVT right-hand side of \eqref{eq:dVdt}--\eqref{eq:complexity}, with additive diffusion and
the measurement model of \S\ref{sec:emp-measurement} as the observation
equation:
\begin{align*}
\dd V &= \Bigl[R\Bigl(1-\tfrac{V}{\Vmax}\Bigr)
        -\delta\,O^{\gamma}\tfrac{V}{V+\varepsilon}\Bigr]\dd t
        + \sigma_V\,\dd W_V, \\
\dd C &= \bigl[\alpha I - \mu C\bigr]\dd t + \sigma_C\,\dd W_C,
\qquad O=O_0+\beta C^{\eta},\quad I=\tfrac{CV}{1+\varphi O}.
\end{align*}
The stochastic extension is a declared departure from the deterministic
theory (\S\ref{sec:limitations}); the diffusion terms absorb unmodeled
shocks, and the drift carries all theoretical content.

\paragraph{Target the dimensionless groups, never the raw parameters.}
The eleven raw parameters are not jointly identifiable from trajectories:
as shown in \S\ref{sec:limitations}, they collapse to approximately six
effective dimensionless groups after nondimensionalizing $V$ by $\Vmax$,
$C$ by $\beta^{-1/\eta}$, and time by $\mu^{-1}$, with pronounced likelihood
ridges along $\mu/\alpha$ (which pins $V^{*}$ via \eqref{eq:vstar}) and
$R/\delta$ (which pins the flux threshold of
Proposition~\ref{prop:carrying}). Estimation therefore parameterizes the
model directly in the reduced groups (as computed by the companion
package's \texttt{dlvt.nondimensional} module), with the raw vector
recovered, when needed, only up to the declared equivalence classes.
Reporting raw-parameter point estimates without this reduction would
manufacture false precision.

\paragraph{Parameter recovery as a mandatory prerequisite.}
Before any field data are collected, the full estimation pipeline must be
validated on synthetic data generated from the model at the proposed
design: the planned number of participants, cohorts, waves, burst lengths,
jittered observation times, measurement reliabilities, and realistic
missingness. The companion software ships an executable
parameter-recovery demonstration (\texttt{dlvt/recovery.py} in the
repository) that simulates trajectories, fits the state-space model, and
reports recovery of the reduced groups and the posterior-correlation
structure. The pre-registered rule is simple: \emph{if the targeted groups
cannot be recovered from data the design would produce, the design---not
the analysis---is revised}, and no field wave is launched until recovery
succeeds. Precedents for fitting dynamic-systems models to intensive
psychological data exist \citep{cramer2016depression,veldhuis2020burnout};
the recovery gate is what separates a fit from an estimate.

\subsection{Identification: breaking the ridges with exogenous variation}
\label{sec:emp-identification}

Cross-sectional equilibrium data identify $\mu/\alpha$ and $R/\delta$ but
not their components (\S\ref{sec:limitations}); trajectories help but leave
the ridges nearly flat. Three sources of exogenous variation break them.

\begin{enumerate}
\item \emph{A status-capital shock identifies $\mu$.} Reorganizations,
mergers and acquisitions, and---cleanest of all---title-only promotions
(status raised with duties and assessed competence unchanged) perturb $C$
without operating through the accumulation channel $\alpha I$. The
post-shock relaxation of the capital indicators toward the nullcline
identifies the depreciation rate $\mu$ directly, and $\alpha$ then follows
from the ratio. Event timing must be plausibly exogenous to the leader's
own $V$ trajectory; firm-level events (reorg announced above the focal
leader's level) are preferred to self-initiated moves.
\item \emph{A randomized recovery manipulation identifies $R$.} A
sabbatical or structured detachment program
\citep{davidson2010sabbatical,sonnentag2022recovery}, randomly assigned,
enters the vitality equation only: it shifts $R$ (and possibly the
diffusion of $V$) while appearing nowhere in the capital equation except
through $V$ itself. This cross-equation exclusion restriction separates
$R$ from $\delta$, exactly as the effort--recovery mechanism
\citep{meijman1998effort} requires, and doubles as the $R\uparrow$ arm of
the killer test (\S\ref{sec:emp-killer}).
\item \emph{An instrument for impact.} $I$ appears in the capital equation
and is measured with error correlated with rater knowledge of the
leader's visible energy, threatening both attenuation and reverse-causal
contamination. Exogenous demand shifters at the unit level---market or
regulatory shocks to the unit's addressable performance, peer-unit
average performance in the same firm---serve as instruments: they move
objective unit outcomes without operating through the focal leader's
current $V$.
\end{enumerate}

\paragraph{Posterior-correlation audit.}
The identification plan carries a quantitative pass/fail criterion: in the
fitted Bayesian model, the posterior correlation between $\mu$ and $\alpha$
must satisfy $|\rho(\mu,\alpha)|<0.6$ (and analogously for $R,\delta$).
If the audit fails, the honest report is the ridge---estimates of
$\mu/\alpha$ and $R/\delta$ only---and any conclusion requiring the
components separately (e.g.\ P6's intervention ordering between
$\mu$-side and $\alpha$-side levers) is withheld.

\subsection{The killer test: a pre-registered $2\times2\times$time field
experiment}
\label{sec:emp-killer}

The decisive confrontation is between scope absorption
(Corollary~\ref{cor:scope}; P3) and intervention asymmetry (P6) on one
side and the natural managerial instinct on the other. The design is a
pre-registered $2\times2$ factorial with longitudinal follow-up
(the full draft pre-registration is \texttt{prereg/PREREGISTRATION.md}):

\begin{description}
\item[Arms.] (i) \emph{Scope reduction} ($\beta\downarrow$): structural
delegation of units, interface pruning, span reduction---interventions
that lower the coupling of the leader's capital to complexity load.
(ii) \emph{Recovery support} ($R\uparrow$): protected restoration
capacity---enforced detachment blocks, recovery training, sabbatical
scheduling \citep{sonnentag2007recovery,hahn2011recovery,
davidson2010sabbatical}.
(iii) \emph{Both}. (iv) \emph{Waitlist control}, receiving the recovery
program after the endpoint.
\item[Endpoint.] Primary endpoint at $12$--$18$ months post-randomization.
The horizon is dictated by the re-expansion mechanism: the transient bump
from $\beta\downarrow$ decays only as capital re-expands, which takes two
to three capital relaxation times ($1/\mu$) plus a substantial fraction of
the spiral period ($\approx19$ time units, about $4.4$ months under the
illustrative weekly calibration). Twelve to eighteen months covers at
least $2$--$3$ capital relaxation times even if the true timescales run
several times slower than the illustrative calibration; a shorter
endpoint would sample the transient and mistake it for a durable effect.
\item[Predictions.] DLVT predicts: (a) a durable rise in vitality in the
$R\uparrow$ arms at the endpoint; (b) in the $\beta\downarrow$-only arm, a
\emph{transient} vitality bump (visible at roughly $3$ months) that
decays as capital re-expands and the product $\beta C^{\eta}$
re-asserts itself; (c) conservation of $\beta C^{\eta}$ in the
$\beta\downarrow$-only arm---measured complexity $O$ returns to its
pre-treatment level as capital indicators grow, per
Corollary~\ref{cor:scope} ($\beta C^{*}=8.008$ at baseline,
$\varepsilon=0.1$).
\item[Informative null.] Because prediction (b) is a null at the endpoint,
the analysis pre-declares a smallest effect size of interest (SESOI) and
tests \emph{equivalence} \citep{lakens2017equivalence} of the
$\beta\downarrow$-only arm to control at
$12$--$18$ months, not merely non-significance. A null that falls inside
the declared equivalence bounds supports P3; an uninstrumented
non-significant difference would support nothing.
\item[Power.] Power is established by simulation from the fitted
state-space model (never by analytic shortcuts that ignore the dynamics):
synthetic trials at the proposed design are generated under (i) the DLVT
drift and (ii) a rival ``durable scope effect'' drift, and the design must
detect the $R\uparrow$ effect, affirm equivalence for $\beta\downarrow$,
and detect the $3$-month transient, each with adequate probability. The
working sketch is $60$--$100$ leaders per arm followed over at least $8$
quarterly waves; the exact numbers are outputs of the simulation, fixed
in the pre-registration before enrollment.
\end{description}

\paragraph{What kills the theory.}
The theory-killing outcome is explicit: \emph{a durable vitality gain in
the scope-reduction-only arm at $12$--$18$ months, with the product
$\beta C^{\eta}$ not conserved}---capital not re-expanding, complexity
durably lowered, vitality durably raised. That outcome falsifies the
re-expansion mechanism itself (P3 and P6 jointly), not a calibration
choice, and no kernel-family retreat (P4) rescues it, because the
conservation claim is the signature of the mechanism rather than of the
power-law family alone. Conversely, a durable scope effect \emph{with}
$\beta C^{\eta}$ conserved would indicate a mis-specified kernel (P4
territory), and a transient-then-decay pattern with conservation is the
theory's fingerprint.

\subsection{Threats to validity, attrition as data, and ethics}
\label{sec:emp-threats}

\paragraph{Reverse causality.}
The theory itself asserts bidirectional coupling ($V$ gates $I$, which
builds $C$, which loads $O$, which drains $V$), so directionality cannot be
settled by lag ordering alone. The defenses are structural: continuous-time
estimation of the full cross-coupled drift rather than interval-dependent
cross-lagged coefficients, plus the exogenous shocks and randomized
manipulation of \S\ref{sec:emp-identification}, which move one equation at
a time.

\paragraph{Survivorship and informative attrition.}
Leaders at low vitality are more likely to exit---turnover, health leave,
demotion, or managed departure---so a naive panel over-samples resilient
survivors and biases $V^{*}$ upward
\citep{conference2023tenure,tu2025dancing,muller2026systematic}. Attrition
is therefore modeled, not deleted: a competing-risks survival sub-model
(voluntary exit, involuntary exit, internal move out of scope) is estimated
jointly with the state-space model, with the hazard depending on the latent
$V$ trajectory (a shared-parameter structure;
\citealp{rizopoulos2012joint}). This turns a threat into a
bonus test: DLVT predicts exit hazard concentrated where the depletion
ratio \eqref{eq:gamma-ratio} exceeds one and during the transient trough of
the damped spiral, and a hazard profile unrelated to latent vitality would
itself be evidence against the model's phenomenology.

\paragraph{Cadence aliasing.}
A rigid quarterly grid samples the illustrative $19$-unit spiral below its
Nyquist rate and could alias oscillation into spurious slow drift. The
mitigations of \S\ref{sec:emp-design}---jittered wave timing, embedded
daily bursts, continuous-time likelihoods over exact intervals, and
pilot-based cadence revision---are pre-committed, and the aliasing risk is
re-assessed once pilot estimates of the empirical timescales exist.

\paragraph{Construct validity.}
The principal exposures are the scalar-$C$ conflation (addressed by the
pre-registered bifactor rival, \S\ref{sec:emp-measurement}), the
discriminability of vitality from reversed exhaustion
\citep{maslach1981burnout}, and the gap between the structural coupling
$\beta$ and its operationalization through complexity indicators; demand
heterogeneity \citep{podsakoff2023challenge,ong2023configural,bakker2017jdr,
demerouti2001job} is recorded so that kernel mis-specification (P4) can be
distinguished from measurement failure.

\paragraph{Ethics and governance.}
The program studies people at their most professionally exposed, and the
governance is part of the design, not an afterthought. (i) A strict
research--HR firewall: individual-level vitality, exhaustion, and
survival-model outputs are never available to the employer and never enter
personnel decisions; only pre-agreed aggregate reports leave the research
environment. (ii) Opt-in consent with per-wave re-consent, and separate
consent for passive metadata (calendar and network indicators), with data
minimization and on-site aggregation. (iii) Neutral participant-facing
language throughout: the construct is described as workload, recovery, and
sustainable performance; theoretical labels such as ``low-vitality
equilibrium''---and, a fortiori, the retired legacy label of earlier
drafts (Definition~\ref{def:lowvitality})---never appear in participant
materials. (iv) The control condition is a waitlist: control-arm
participants receive the recovery program after the primary endpoint, so no
one is denied a plausibly beneficial intervention, merely scheduled later.
(v) Institutional review and, where applicable, works-council approval
precede enrollment, and the pre-registration binds the analysis before the
first participant is randomized.
